\chapter{Graphical Models and Message Passing}
\label{ch:graphs}\label{ch:graphical}

This chapter introduces the representation the thesis computes in: a
distribution written as a product of local factors, and the recursion that
returns its marginals exactly when those factors form a tree. It is short by
design. The multivariate Gaussian in both parametrisations, the spectral and
matrix-inversion tools, the Gaussian special case of the recursion, and its
relation to TAP and AMP are all in Appendix~\ref{app:linalg}; what is kept here
is the structure, because it is the structure and not the algebra that decides
what is computable.

\section{Markov random fields}\label{sec:mrf}

Appendix~\ref{app:linalg} assigns structure to matrices, through the zeros of
a precision. Graphical models assign it to distributions instead, for arbitrary
and not necessarily Gaussian variables; the two turn out to be the same
assignment.

\begin{definition}[Markov random field]
Let $G = (V, \mathcal{E})$ be an undirected graph with one random
variable $x_i$ per node. The collection $x = (x_i)_{i \in V}$ is a
\emph{Markov random field} (MRF) if it satisfies the local Markov
property: for every node $i$,
\begin{equation}
  p\big(x_i \,\big|\, x_{V \setminus \{i\}}\big)
  \;=\; p\big(x_i \,\big|\, x_{\partial i}\big),
  \label{eq:local-markov}
\end{equation}
where $\partial i$ is the set of neighbours of $i$ in $G$.
\end{definition}

The fundamental structure theorem connects this
conditional-independence definition to the energy picture of
the energy picture: by the Hammersley--Clifford theorem
\citep[proved in][]{besag1974spatial}, a strictly positive density is an
MRF on $G$ if and only if it factorises over the \emph{cliques} of $G$,
\begin{equation}
  p(x) \;=\; \frac{1}{Z} \prod_{C \in \mathrm{cliques}(G)} \psi_C(x_C)
  \;=\; \frac{1}{Z}\, \exp\Big[ -\sum_C E_C(x_C) \Big] .
  \label{eq:hammersley-clifford}
\end{equation}
Conditional independence, graphical locality, and additive local
energies are one and the same thing. The Ising model \eqref{eq:ising} is
the textbook example. The decisive example for us is one-dimensional:
the chain factorisation \eqref{eq:chain-factorisation} makes a Markov
chain an MRF on a line, with pairwise nearest-neighbour energies. For
\emph{Gaussian} MRFs the correspondence becomes matrix-concrete: the
energy is the quadratic form $\tfrac12 a^\T J a - h^\T a$, its pairwise
terms are the entries $J_{ij}$, and the graph of the MRF is exactly the
sparsity pattern of the precision matrix. The dual reading of zeros in
Appendix~\ref{sec:gauss-vectors} and the Hammersley--Clifford theorem are
the same statement, once in linear algebra and once in probability. A
Gaussian Markov chain therefore has a tridiagonal precision, a fact
Chapter~\ref{ch:gaussian} derives explicitly and exploits throughout.

One more structural remark that the research chapters rely on: if each
variable of an MRF is observed through its own independent noisy channel,
the likelihood factors touch one variable each, so no clique is
enlarged. \emph{The posterior of an MRF under per-node observations is an
MRF on the same graph.} Observation changes the potentials, never the
graph.

\section{Factor graphs}\label{sec:factor-graphs}

MRFs record \emph{which} variables interact but not \emph{how} the joint
factorises: a triangle of pairwise terms and a single triple term have
the same graph. The representation that makes the factorisation itself
graphical, and on which message passing is most naturally defined, is
the \emph{factor graph} \citep{kschischang2001factor,
koller2009probabilistic}.

\begin{definition}[Factor graph]
Given a factorisation
$p(x) = \frac{1}{Z}\prod_{a} \psi_a(x_{\partial a})$, its factor graph is
the bipartite graph with a \emph{variable node} for each $x_i$, a
\emph{factor node} for each $\psi_a$, and an edge $(i, a)$ whenever
$x_i \in x_{\partial a}$.
\end{definition}

A connected factor graph with one fewer edge than nodes is a
\emph{tree}: it contains no cycles, so any edge separates it into two
disjoint components. Trees are the tractable case of inference: on
them, marginals and partition functions are computable exactly by the
recursion of the next section, whatever the alphabet of the variables.
The factor graph of a Markov chain observed through per-node noise is a
tree of a particularly simple shape (a spine of variables and transition
factors, with one evidence leaf per variable), as
Chapter~\ref{ch:gaussian} verifies by counting.

\section{The sum--product algorithm}\label{sec:sum-product}

On a factor graph, sum--product belief propagation (BP) passes messages
along every edge \citep{pearl1988probabilistic, kschischang2001factor,
mezard2009information}. Variable-to-factor messages multiply the
incoming factor messages from all \emph{other} neighbours, and
factor-to-variable messages integrate the factor against the incoming
variable messages:
\begin{equation}
  \nu_{i \to a}(x_i) \propto
    \prod_{b \in \partial i \setminus a} \hat\nu_{b \to i}(x_i),
  \qquad
  \hat\nu_{a \to i}(x_i) \propto
    \int \psi_a(x_{\partial a})
    \prod_{j \in \partial a \setminus i} \nu_{j \to a}(x_j)\;
    \dd x_{\partial a \setminus i},
  \label{eq:bp-rules}
\end{equation}
with the belief at node $i$ proportional to
$\prod_{a \in \partial i} \hat\nu_{a \to i}(x_i)$. The ``exclude the
recipient'' rule is the heart of the algorithm: the message $i \to a$
summarises everything node $i$ knows \emph{except} what it learned from
$a$, so that no piece of evidence is ever counted twice along a path.

\begin{theorem}[Tree exactness of sum--product]\label{thm:tree-exact}
On a tree-shaped factor graph, the equations \eqref{eq:bp-rules} have a
unique solution, reached after one inward and one outward pass, and the
resulting beliefs are the exact marginals of $p$
\citep[Theorem~14.1]{mezard2009information}.
\end{theorem}

The reason is the separation property of trees: each edge splits the
graph into two disjoint halves, and the message crossing it summarises
exactly the half behind it. On graphs with cycles the same updates,
iterated, define \emph{loopy} BP, which is extraordinarily useful but
approximate \citep{weiss2001correctness, malioutov2006walk}; none of that
is needed in this thesis, because our graphs are trees.

\paragraph{The representation problem.} For \emph{discrete} variables
each message is a finite probability vector and each update in
\eqref{eq:bp-rules} is finite arithmetic. For \emph{continuous} variables
each message is a probability \emph{density}, an infinite-dimensional
object, and the update integrals must be carried out in function space.
Tree exactness (a property of the graph) says nothing about
representability of the messages (a property of the model's function
class); keeping the two apart is essential, and the distinction is where
all the difficulty of the non-Gaussian research chapter lives. There are
three ways to make continuous BP computable: restrict to models whose
messages close in a finite-dimensional family (the Gaussian route, next
section); represent messages numerically on a grid (the route our
non-Gaussian reference solver takes, Chapter~\ref{ch:laplace}); or
project each message onto a tractable family and accept the error (the
route whose cost this thesis measures).

\section{Locality in learned architectures: convolutions, receptive
fields, U-Nets}\label{sec:cnn-locality}

The message-passing story has an architectural mirror, and it supplies
the vocabulary for the locality questions of the research chapters. A
convolutional network processes a signal by applying the same local
filter at every position, stacking such layers so that depth enlarges
the \emph{receptive field}: the window of input coordinates that can
influence a given output coordinate \citep{lecun2015deep}. Convolutional
architectures powered the breakthrough of deep learning in vision
\citep{krizhevsky2012imagenet}, and their central inductive bias is
precisely locality plus translation invariance: the value of an output
pixel is computed from a bounded neighbourhood, identically across the
image. The standard denoising backbone of diffusion models, the U-Net
\citep{ronneberger2015unet, ho2020denoising}, is a convolutional
encoder--decoder with skip connections, whose receptive field grows
along the contracting path and whose skip connections reinject
fine-scale local information along the expanding path.

Why local architectures should be \emph{sufficient} for score estimation
is a theoretical question, and message passing has begun to answer it.
\citet{mei2024unets} shows that for generative hierarchical models
(tree-structured graphical models over image-like data), the denoising
function computed by belief propagation can be implemented efficiently
by a U-Net, with the network's layers mirroring the sweeps of the
message-passing schedule; convolutional layers are efficient
implementations of the local BP updates. This gives the U-Net an
inference-algorithm reading: the architecture is good at denoising
because denoising, on structured data, \emph{is} message passing, and
message passing is local.

The correspondence is not confined to convolutional architectures.
\citet{garnierbrun2024transformers} train transformers on data generated
by a hierarchical model in which the exact posterior is again computable
by belief propagation, and find that the trained attention maps come to
reproduce the message-passing schedule. Read together with
\citet{mei2024unets}, this is \emph{cross-architecture evidence} for a
particular reading: that graphical inference is a useful description of what a
trained denoiser ends up computing, and that the depth it needs is set by the
structure of the data rather than by the architecture chosen. It is evidence
across two architecture families and two model families, not a theorem about
every denoiser --- a network can approximate the same posterior map without
reproducing the algorithm internally, and neither result rules that out. This is
precisely why the chain matters here. A chain is a
degenerate tree, of depth equal to its length and width two, so the depth
these results would predict is not the logarithmic depth of a balanced
hierarchy but a depth linear in the correlation length --- a sharply
different architectural prediction, on a model where it can be checked.

The present thesis works one step below this correspondence, on a model
small enough that the locality of the score is not an architectural
hypothesis but a computable quantity. For the Gaussian chain, the exact
score at frame $k$ depends on the observation at frame $k \pm d$ with a
weight that decays geometrically in $d$, at a rate $q(\alpha, t)$
computed in closed form in Chapter~\ref{ch:gaussian}. The error of the
radius-$r$ windowed estimator is \emph{measured} to decay at that same rate
over the tested range --- a separate and weaker statement. Comparing
truncation of the inference against truncation of the score matrix then
quantifies, in one solvable case, how a local computation should be
organised. These results describe the function a network would have to
learn rather than the dynamics of learning it; the trained baselines of
Chapter~\ref{ch:em-results} supply the other half of that comparison.
